\documentclass[conference,a4paper]{IEEEtran}
\usepackage[T1]{fontenc}
\usepackage[utf8]{inputenc}
\usepackage[brazil]{babel}
\usepackage{microtype}
\usepackage{amsmath,amssymb}
\usepackage{booktabs}
\usepackage{array}
\usepackage{graphicx}
\usepackage{balance}
\usepackage{xcolor}
\usepackage{hyperref}
\hypersetup{colorlinks=true,linkcolor=black,urlcolor=blue!55!black,citecolor=blue!45!black,pdftitle={Atenção em séries temporais: previsão climática com Transformer Encoder},pdfauthor={Luan Diniz Mazaro Rodovalho}}
\renewcommand{\arraystretch}{1.12}
\setlength{\tabcolsep}{3pt}
\title{Atenção em séries temporais: previsão meteorológica multivariada com Transformer Encoder}
\author{\IEEEauthorblockN{Luan Diniz Mazaro Rodovalho}
\IEEEauthorblockA{Universidade Federal de Uberlândia (UFU)\\Brasil}}
\begin{document}
\maketitle
\begin{abstract}
Este trabalho investiga o uso de um codificador Transformer para previsão multivariada de dados meteorológicos da série Jena Climate 2009--2016. São considerados dois objetivos: prever o próximo token de uma hora e prever conjuntamente 24 tokens futuros, correspondentes a 24 horas. As medições são organizadas em janelas temporais, normalizadas com estatísticas exclusivas do treinamento e processadas por blocos de autoatenção. Foram realizados seis experimentos, incluindo variações do histórico, da dimensão de representação, do número de cabeças e da profundidade. Na previsão de uma hora, a configuração com quatro blocos obteve RMSE de temperatura de 0,504~\textdegree C; para 24 horas, o modelo inicial obteve 2,412~\textdegree C. Curvas de perda, exemplos de previsão e mapas médios de atenção são utilizados para discutir aprendizado, dependências temporais e limitações.
\end{abstract}
\begin{IEEEkeywords}
Transformer, séries temporais, autoatenção, previsão meteorológica, aprendizado profundo.
\end{IEEEkeywords}

\section{Introdução}
O objetivo deste mini-projeto é investigar um \emph{Transformer Encoder} aplicado ao conjunto Jena Climate 2009--2016. Em particular, avaliam-se (i) a previsão do próximo token temporal e (ii) a previsão direta de tokens consecutivos que representam 24 horas. Também se examina o efeito de quatro hiperparâmetros: número de tokens de entrada, dimensão do \emph{embedding}, quantidade de cabeças de atenção e número de blocos codificadores. Além dos erros de previsão, são analisadas as curvas de perda e matrizes de atenção, conforme o enunciado do projeto.

A implementação foi desenvolvida em Python/PyTorch e está disponível em \url{https://github.com/LuanDinizMR/projeto-transformers}.

\section{Dados e preparação}
\subsection{Dataset, estatísticas e limpeza}
Utilizou-se o conjunto \emph{Jena Climate 2009--2016}, composto por cerca de 420 mil registros coletados em intervalos de dez minutos e 14 variáveis meteorológicas, como pressão, temperatura, umidade e vento \cite{jena}. A Tabela~\ref{tab:estatisticas} apresenta médias e desvios-padrão de variáveis selecionadas, obtidos das estatísticas do conjunto de treinamento armazenadas com o modelo. Os valores descrevem apenas esse subconjunto, não toda a série.

\begin{table}[!t]
\centering
\caption{Estatísticas do conjunto de treinamento após limpeza.}
\label{tab:estatisticas}
\footnotesize
\begin{tabular}{@{}lrr@{}}
\toprule
Variável & Média & Desvio-padrão \\
\midrule
Pressão (mbar) & 988,656 & 8,297 \\
Temperatura ($^\circ$C) & 9,108 & 8,654 \\
Umidade relativa (\%) & 75,904 & 16,557 \\
Velocidade do vento (m/s) & 2,155 & 1,530 \\
Rajada máxima (m/s) & 3,561 & 2,314 \\
Direção do vento (graus) & 173,442 & 87,674 \\
\bottomrule
\end{tabular}
\end{table}

Na inspeção dos dados, foram identificados 18 valores sentinela $-9999$ na velocidade do vento e 20 na velocidade máxima do vento, todos concentrados no intervalo destinado à validação. Esses valores representam medições inválidas e foram substituídos por zero, seguindo uma convenção utilizada para este conjunto de dados. Sem a limpeza, a perda de validação permanecia próxima de 1.295, apesar de a perda de treinamento ser inferior a 0,12; após a correção, ambas passaram a apresentar escalas comparáveis. Esse diagnóstico reforça a importância do pré-processamento antes da avaliação do modelo.

\subsection{Partição e tokenização temporal}
A série foi dividida cronologicamente em 70\% para treinamento, 15\% para validação e 15\% para teste. Para cada variável $j$, a transformação $z=(x-\mu_j)/\sigma_j$ foi ajustada \emph{somente} nos dados de treinamento e aplicada aos três conjuntos. Não houve embaralhamento dos registros antes da separação. As janelas são construídas dentro de cada partição, sem cruzar suas fronteiras.

Cada token contém seis medições consecutivas (uma hora), com 14 variáveis por medição. O vetor bruto de token tem, portanto, $6\times14=84$ atributos. Não foi empregada média horária: os seis registros são concatenados para preservar a granularidade de dez minutos. A configuração inicial utiliza $m=24$ tokens (histórico $L=24$ horas); outra configuração utiliza $m=48$ ($L=48$ horas). Janelas deslizantes avançam um token de cada vez, produzindo 49.040 exemplos de treino, 10.489 de validação e 10.489 de teste para o objetivo de uma hora. No objetivo de 24 horas, cada saída contém 24 tokens, isto é, $24\times84=2.016$ valores.

\section{Modelo e procedimento experimental}
\subsection{Arquitetura do codificador}
Cada token de entrada $\mathbf{x}_t\in\mathbb{R}^{84}$ é transformado por projeção linear em $\mathbf{e}_t\in\mathbb{R}^{d}$, à qual se soma uma codificação posicional senoidal. Em seguida, os $m$ tokens são processados por blocos codificadores. A autoatenção de cada cabeça utiliza consultas, chaves e valores produzidos a partir das representações dos tokens:
\begin{equation}
 \mathrm{Attention}(Q,K,V)=\mathrm{softmax}\!\left(\frac{QK^\top}{\sqrt{d_k}}\right)V.
\end{equation}
As várias cabeças permitem construir diferentes representações de relacionamentos temporais. Cada bloco emprega \emph{LayerNorm} antes da atenção e da rede \emph{feed-forward}, conexões residuais, \emph{dropout} de 0,1 e uma rede interna com expansão de $d$ para $4d$, ativação GELU e retorno a $d$. Trata-se de autoatenção bidirecional \emph{dentro da janela histórica}; nenhum token futuro é fornecido ao codificador.

Após os blocos, utiliza-se a representação do último token histórico, seguida de normalização e de uma camada linear que prevê diretamente $h\times84$ valores, reorganizados em $h$ tokens futuros. Assim, $h=1$ e $h=24$ são resolvidos pelo mesmo desenho de rede, mas com diferentes dimensões de saída. A previsão de 24 horas é conjunta, sem realimentação autoregressiva dos valores previstos.

\subsection{Treinamento, hiperparâmetros e métricas}
Todos os treinamentos usam otimizador Adam, taxa inicial de aprendizado $10^{-3}$, lote de 64 exemplos, 20 épocas e semente pseudoaleatória 42. A execução informada utilizou CPU. A função de perda é o erro quadrático médio (MSE), adequado à regressão de variáveis contínuas e calculado sobre os valores normalizados:
\begin{equation}
 \mathrm{MSE}=\frac{1}{N}\sum_{i=1}^{N}(\hat y_i-y_i)^2.
\end{equation}
O MSE penaliza mais intensamente erros grandes. Os pesos selecionados são os da época com menor perda de validação. A avaliação final usa RMSE e MAE de todas as variáveis na escala normalizada; para temperatura, os valores previstos e reais são desnormalizados para reportar RMSE e MAE em graus Celsius.

A Tabela~\ref{tab:config} mostra as seis configurações testadas. Em E3--E6 varia-se um hiperparâmetro por vez em relação a E1, mantendo o horizonte em uma hora; E2 compara outro horizonte com a arquitetura inicial.

\begin{table}[!t]
\centering
\caption{Configurações experimentais. $m$: tokens históricos; $h$: tokens futuros; $d$: dimensão; $c$: cabeças; $b$: blocos.}
\label{tab:config}
\footnotesize
\begin{tabular}{@{}llrrrrr@{}}
\toprule
ID & Descrição & $m$ & $h$ & $d$ & $c$ & $b$ \\
\midrule
E1 & Inicial (1h) & 24 & 1 & 64 & 4 & 2 \\
E2 & Horizonte de 24h & 24 & 24 & 64 & 4 & 2 \\
E3 & Mais histórico & 48 & 1 & 64 & 4 & 2 \\
E4 & Maior dimensão & 24 & 1 & 128 & 4 & 2 \\
E5 & Mais cabeças & 24 & 1 & 64 & 8 & 2 \\
E6 & Mais blocos & 24 & 1 & 64 & 4 & 4 \\
\bottomrule
\end{tabular}
\end{table}

\section{Resultados e discussão}
\subsection{Erros de previsão}
Os resultados do conjunto de teste são mostrados na Tabela~\ref{tab:metricas}. Na previsão de uma hora, E6 apresentou os menores erros \emph{de temperatura}, com RMSE de 0,504~$^\circ$C e MAE de 0,343~$^\circ$C, reduções de aproximadamente 2,5\% e 3,4\% em relação ao E1. Entretanto, E1 obteve o menor RMSE normalizado considerando conjuntamente as 14 variáveis (0,256356, contra 0,256738 de E6). Não há, portanto, superioridade uniforme de uma configuração em todas as métricas.

\begin{table}[!t]
\centering
\caption{Resultados no teste. As duas primeiras métricas abrangem todas as variáveis em escala normalizada; as duas últimas se referem à temperatura.}
\label{tab:metricas}
\footnotesize
\begin{tabular}{@{}lrrrr@{}}
\toprule
 & \multicolumn{2}{c}{Todas as variáveis} & \multicolumn{2}{c}{Temperatura ($^\circ$C)} \\
ID & RMSE & MAE & RMSE & MAE \\
\midrule
E1 & \textbf{0,256356} & 0,113388 & 0,517569 & 0,354476 \\
E2 & 0,506768 & 0,330482 & 2,411583 & 1,825502 \\
E3 & 0,257965 & 0,116573 & 0,562217 & 0,401988 \\
E4 & 0,256868 & 0,114259 & 0,522154 & 0,356324 \\
E5 & 0,257362 & 0,113841 & 0,520744 & 0,357784 \\
E6 & 0,256738 & \textbf{0,112493} & \textbf{0,504389} & \textbf{0,342578} \\
\bottomrule
\end{tabular}
\end{table}

Ampliar o histórico para 48 horas (E3) aumentou o erro de temperatura em relação a E1; aumentar a dimensão para 128 (E4) ou o número de cabeças para oito (E5) também não trouxe melhora nessas métricas. Já a previsão direta de 24 horas (E2) apresentou RMSE de temperatura de 2,412~$^\circ$C, aproximadamente 4,7 vezes o valor de E1. A tarefa é distinta: nela o modelo precisa gerar, de uma só vez, uma sequência futura muito mais extensa, de modo que a comparação não constitui uma avaliação isolada de hiperparâmetros.

\subsection{Curvas de aprendizagem}
A Figura~\ref{fig:perdas} apresenta as curvas de MSE do modelo inicial de uma hora e da previsão de 24 horas. Em E1, a perda de treino diminui rapidamente nas primeiras épocas e depois se estabiliza; a validação oscila em nível próximo, sem a discrepância extrema observada antes da correção dos valores sentinela. Em E2, embora a perda de treino siga diminuindo, a curva de validação apresenta seu mínimo nas primeiras épocas e tende a crescer posteriormente. Esse afastamento é compatível com \emph{overfitting}, motivo pelo qual foi utilizado o \emph{checkpoint} com melhor validação.

\begin{figure*}[!t]
\centering
\includegraphics[width=0.93\textwidth]{figuras/perdas_comparado.png}
\caption{Perdas de treinamento e validação em duas tarefas. Os gráficos originais usam índices de época começando em zero.}
\label{fig:perdas}
\end{figure*}

A Figura~\ref{fig:previsoes} mostra exemplos individuais de temperatura. No horizonte curto, as previsões ficam próximas dos registros reais na hora seguinte. Em 24 horas, o modelo acompanha parte da tendência geral, mas suaviza oscilações e apresenta divergências de magnitude e momento. Essas visualizações são exemplos ilustrativos do comportamento dos modelos, enquanto as métricas da Tabela~\ref{tab:metricas} consideram todo o conjunto de teste.

\begin{figure*}[!t]
\centering
\includegraphics[width=0.93\textwidth]{figuras/temperatura_prevista_comparado.png}
\caption{Temperatura real e prevista em uma amostra de teste para cada horizonte.}
\label{fig:previsoes}
\end{figure*}

\subsection{Mapas médios de atenção}
Foram extraídos pesos de atenção da \emph{primeira cabeça do último bloco} para as janelas do conjunto de teste. Em seguida, realizou-se a média dessas matrizes sobre todas as amostras, obtendo mapas $m\times m$, com as posições de consulta nas linhas e de chave nas colunas. A Figura~\ref{fig:atencao} compara os dois horizontes com a configuração inicial de arquitetura. No E1, observa-se concentração de pesos nas posições finais da janela, sugerindo que essa cabeça recorre intensamente às horas mais recentes. No E2, a atenção está mais distribuída pelo histórico, mantendo maior peso em regiões próximas do fim da janela.

\begin{figure}[!t]
\centering
\includegraphics[width=\columnwidth]{figuras/atencao_horizontes.png}
\caption{Mapas de atenção médios para E1 (1h) e E2 (24h); cada painel corresponde à primeira cabeça da última camada, com sua própria escala de cores.}
\label{fig:atencao}
\end{figure}

A Figura~\ref{fig:cabecas} evidencia que alterar o número de cabeças também modifica a distribuição observada: no E5, a primeira cabeça concentra pesos em posições mais antigas, sobretudo perto do quarto token, enquanto no E1 a mesma posição de cabeça privilegia as horas finais. Isso não indica que uma distribuição seja intrinsecamente melhor, como confirmam os erros de temperatura próximos entre E1 e E5.

\begin{figure}[!t]
\centering
\includegraphics[width=\columnwidth]{figuras/atencao_cabecas.png}
\caption{Primeira cabeça da última camada nos modelos com quatro (E1) e oito (E5) cabeças, ambos prevendo uma hora.}
\label{fig:cabecas}
\end{figure}

Os mapas não são explicações causais das previsões: representam apenas uma das cabeças e uma agregação sobre muitas janelas. Cabeças diferentes podem apresentar outros padrões. Além disso, a média tende a ocultar comportamentos particulares de eventos meteorológicos e horários específicos. Ainda assim, os mapas demonstram que os pesos de autoatenção podem ser recuperados e associados à organização temporal das entradas.

\balance
\section{Conclusão}
Foi implementado e avaliado um Transformer Encoder para duas tarefas de previsão meteorológica multivariada sobre janelas horárias. A arquitetura combina projeção linear, posição senoidal, autoatenção multi-cabeças, blocos \emph{feed-forward} e camada linear de regressão. Após corrigir 38 valores inválidos de vento concentrados na validação, as curvas de perda se tornaram coerentes, permitindo interpretar os treinamentos. Nas configurações de uma hora, o modelo com quatro blocos apresentou o menor RMSE de temperatura (0,504~$^\circ$C); entretanto, o modelo inicial foi ligeiramente melhor no RMSE normalizado global. No horizonte de 24 horas, os erros foram substancialmente maiores e surgiram sinais de sobreajuste nas curvas de validação.

O estudo é limitado por uma única semente de execução, pela ausência de intervalos de confiança e pela apresentação de apenas uma cabeça de atenção por configuração. Como continuação, seria possível avaliar diferentes sementes, quantificar incerteza e comparar outras estratégias de previsão; tais extensões não são necessárias para os objetivos experimentais aqui apresentados.

\begin{thebibliography}{9}
\bibitem{vaswani2017} A. Vaswani et al., ``Attention Is All You Need'', \emph{Advances in Neural Information Processing Systems}, 2017. \url{https://arxiv.org/abs/1706.03762}.
\bibitem{jena} TensorFlow, ``Jena Climate 2009--2016'' (arquivo de dados). \url{https://storage.googleapis.com/tensorflow/tf-keras-datasets/jena_climate_2009_2016.csv.zip}.
\bibitem{pytorch} PyTorch, ``MultiheadAttention'', documentação da biblioteca. \url{https://pytorch.org/docs/stable/generated/torch.nn.MultiheadAttention.html}.
\end{thebibliography}
\end{document}
